\documentclass[11pt,reqno]{amsart}

\usepackage[letterpaper,margin=1in]{geometry}
\usepackage{amsmath,amssymb,mathtools,microtype}
\usepackage[hidelinks]{hyperref}
\hypersetup{pdftitle={Erdos 123+ : distinct 2^k 3^l 5^m summands within a factor 1+epsilon}}

\setlength{\parindent}{1.15em}

\newtheorem{theorem}{Theorem}
\newtheorem{corollary}{Corollary}

\newcommand{\Sg}{\mathcal S}
\newcommand{\N}{\mathbb N}
\newcommand{\Z}{\mathbb Z}
\newcommand{\Nzero}{\mathbb Z_{\ge 0}}

\title[Erd\H{o}s 123$^{+}$]{Erd\H{o}s 123$^{+}$: distinct $2^{k}3^{\ell}5^{m}$
summands within a factor $1+\varepsilon$}
\author{Principia Math}
\date{}

\begin{document}

\begin{abstract}
We record a one-line corollary of the representation theorem behind our solution of
Erd\H{o}s Problem~123: for $\{2^{k}3^{\ell}5^{m}\}$ and every $\varepsilon>0$, every
sufficiently large integer is a sum of distinct such numbers whose largest and smallest
terms differ by a factor less than $1+\varepsilon$.  This settles the stronger conjecture
Erd\H{o}s recorded in~1992 (\cite{Er92b}; the comment to erdosproblems.com/123), by
applying that theorem with $\rho=1+\varepsilon$.
\end{abstract}

\maketitle

\section*{The theorem we proved}

Let $a,b,c\ge2$ be pairwise coprime and
$\Sg=\{a^{k}b^{\ell}c^{m}:k,\ell,m\in\Nzero\}$; write $d=\min(a,b,c)$.  Our solution of
Erd\H{o}s~123 rests on the following theorem, formalized and machine-checked in Lean~4
(Theorems~1.1 and~1.2 of \emph{Antichain Subset Sums in Rank-Three Multiplicative
Semigroups}; for rational $\rho$ it comes from a uniform local central limit theorem for
the Bernoulli subset sums of $\Sg\cap[x,\rho x)$).

\begin{theorem}\label{thm:rep}
For every real $\rho$ with $1<\rho<d$, every sufficiently large integer $N$ can be written
as
\[
  N=\sum_{s\in A}s,\qquad A\subset \Sg\cap[x,\rho x)
\]
for some positive integer $x$.  Since any two elements of $[x,\rho x)$ have ratio
$<\rho<d$, no element of $A$ divides another; thus the terms of $A$ are distinct and no
one of them divides another.
\end{theorem}

\section*{The strengthening, and how it follows}

For $(a,b,c)=(2,3,5)$, Erd\H{o}s asked for more than distinctness and non-divisibility: he
conjectured that all the summands can be taken from a single interval $[x,\rho x)$ with
$\rho$ as close to~$1$ as one wishes.

\begin{corollary}[Erd\H{o}s 123$^{+}$, \cite{Er92b}, 1992]\label{cor:plus}
For every $\varepsilon>0$, every sufficiently large integer $n$ can be written as a sum of
distinct integers $b_{1}<\dots<b_{t}$ of the form $2^{k}3^{\ell}5^{m}$ with
\[
  b_{t}<(1+\varepsilon)\,b_{1}.
\]
\end{corollary}

\begin{proof}
Here $d=\min(2,3,5)=2$.  Given $\varepsilon>0$, set
\[
  \rho=1+\min\!\left(\varepsilon,\tfrac12\right)\in(1,2),
\]
so that $1<\rho<d$ and Theorem~\ref{thm:rep} applies.  For all large $n$ it gives a set
$A\subset\Sg\cap[x,\rho x)$ with $\sum_{s\in A}s=n$; list its distinct elements as
$b_{1}<\dots<b_{t}$.  Every $b_{i}$ lies in $[x,\rho x)$, so $b_{1}\ge x$ and
\[
  b_{t}<\rho x\le(1+\varepsilon)\,x\le(1+\varepsilon)\,b_{1}.
\]
The terms are distinct and pairwise non-dividing by the last sentence of
Theorem~\ref{thm:rep}.  (When $\varepsilon\ge1$ the choice $\rho=\tfrac32$ already gives
the claim.)
\end{proof}

The only freedom used is the choice of $\rho$: the two endpoints of the interval
$[x,\rho x)$ in Theorem~\ref{thm:rep} are in the ratio $\rho$, which we take to be
$1+\varepsilon$.  The same argument gives the corresponding statement for every
pairwise-coprime triple $a,b,c$, since a ratio $\rho=1+\varepsilon<d$ is available whenever
$\varepsilon<d-1$.

\begin{thebibliography}{9}
\bibitem{Er92b}
P.~Erd\H{o}s, conjecture recorded in~1992; see the comments to Erd\H{o}s Problem~123 at
\texttt{https://www.erdosproblems.com/123} (reference \texttt{[Er92b]}).
\end{thebibliography}

\end{document}
